\section{Compresión de datos}

Aunque los métodos de neural SBI eliminan la necesidad de elegir a mano cosas como el umbral $\epsilon$ de ABC, todavía queda un problema importante: los datos $x$ pueden ser de muy alta dimensión, mientras que los parámetros $\Theta$ suelen ser de baja dimensión. Entonces sigue siendo difícil mapear directamente desde $x$ hacia $\Theta$. Igual que en ABC necesitábamos buenas \textit{summary statistics}, acá también necesitamos alguna forma de comprimir los datos. La idea es construir una función $s(x)$ que reduzca la dimensión:
\begin{equation}
x \in \mathbb{R}^{D_x} \rightarrow s(x) \in \mathbb{R}^{D_s}, \quad D_s < D_x
\end{equation}

pero que conserve la mayor cantidad posible de información relevante sobre los parámetros. Idealmente, queremos que trabajar con los datos comprimidos sea casi equivalente a usar los datos originales:
\begin{equation}
p(\Theta | s(x)) \approx p(\Theta | x)
\end{equation}

O sea, que la compresión no destruya información útil para inferencia. En general, hay tres formas principales de abordar esto: usar estadísticas diseñadas a mano basadas en conocimiento del problema, usar métodos que maximizan información como Fisher, o aprender la compresión directamente con redes neuronales junto con el modelo de inferencia.

\subsection{Summary Statistics}

La forma más simple de reducir la dimensión de los datos es usar \textit{summary statistics} diseñadas a mano antes de aplicar neural SBI. La idea es aprovechar conocimiento del dominio para construir funciones $s(x)$ que capturen lo importante. En algunos casos ideales, estas estadísticas pueden ser suficientes, es decir, contienen toda la información necesaria para inferir los parámetros. Pero en problemas reales, esto rara vez pasa. Entonces hay un trade-off claro: las estadísticas hechas a mano suelen ser interpretables, pero pueden perder información importante si no capturan todo lo relevante.

\subsection{Compresión de información teórica}

Para datos más complejos, hay enfoques que en vez de definir las estadísticas a mano, las aprenden maximizando información. La idea general es encontrar una compresión que retenga la mayor cantidad posible de información sobre los parámetros, por ejemplo, usando información de Fisher. Un ejemplo es el enfoque de \cite{heavens2000}, que hace compresión lineal óptima bajo ciertos supuestos y reduce los datos a tantas dimensiones como parámetros haya. Básicamente proyecta los datos en direcciones que maximizan la información. Otro enfoque es el que usan en \cite{alsing2018}, donde primero se extraen features y luego se comprimen hacia el \textit{score}, que es el gradiente del log-likelihood respecto a los parámetros. Esto permite capturar información relevante incluso en casos no gaussianos. También existen métodos como el de \cite{charnock2018}, donde directamente se entrena una red neuronal para maximizar el determinante de la matriz de información de Fisher. Acá el trade-off es distinto: se tiene más flexibilidad y potencialmente mejores representaciones, pero se pierde algo de interpretabilidad comparado con estadísticas hechas a mano.

\subsection{Compresión End-to-End}

Una idea aún más moderna es no separar compresión e inferencia, sino entrenarlas juntas. En vez de definir $s(x)$ antes, se usa una red $s_{\Psi}(x)$ que aprende una representación, y luego otra red $q_{\Phi}$ que hace la inferencia:
\[
x \rightarrow s_{\Psi}(x) \rightarrow q_{\Phi}(\Theta|x)
\]

Ambas redes se entrenan al mismo tiempo usando la misma loss, por ejemplo la de NPE, así que la compresión se adapta automáticamente a lo que necesita la inferencia. Dependiendo del tipo de datos, se pueden usar distintas arquitecturas: CNNs para imágenes, RNNs para secuencias, etc. Este enfoque también se usa en NRE sin problemas. Un detalle importante es que esto no funciona tan bien para NLE: como la likelihood no fuerza a mantener información relevante de los parámetros, el modelo puede colapsar a representaciones poco informativas. Por eso, en NLE suele ser mejor fijar la compresión desde antes. En teoría, este enfoque puede aprender estadísticas casi óptimas, pero en la práctica tiene costos: requiere muchos datos, es menos interpretable, y puede fallar si el modelo se complejiza demasiado con estos pesos y bias extras.


